\subsection{All acyclic categories through twenty-one arrows}

\begin{lemma}[All degree-two products on the two small five-chain levels]
\label{C:five-chain-nineteen-products}
Let $D$ have five linearly ordered objects, trivial endomorphisms,
and a nonempty hom-set in every forward direction.  If $D$ has
at most nineteen arrows, then every product of two classes in
$H^2(BD;\mathbb F_2)$ is zero.  No assumption on its first homology
or the dimension of its second homology is required.
\end{lemma}
\begin{proof}
The top-face argument in Proposition~\ref{C:five-object-acyclic21}
shows that $H_4(BD;\mathbb F_2)=0$ unless
\[
 h_{01},h_{34}\geq2,\qquad h_{12}>1\text{ or }h_{23}>1.
\]
These conditions already require eighteen arrows.  At totals
eighteen and nineteen there are respectively two and nineteen
positive cardinality profiles satisfying them.

The generator \path{code/d_five_object_line_all.cpp} enumerates
all completed tables on these profiles, with no homology filter.
It gives two and sixty-seven labelled associative tables, or two
and thirty-one orbits under permutations within hom-sets.
The standalone checker \path{code/d_five_chain_all_cups.py}
checks every typed composition law and every bar differential,
constructs a full basis of second mod-two cohomology, and finds
an explicit degree-three primitive for every ordered product
of basis classes.  Thus it verifies all products, not just the
square in a previously selected one-dimensional case.

There is also an independent enumeration with no associativity
propagation: \path{code/d_five_chain_independent.py} forms every
Cartesian product of the allowed values of the unprescribed
typed products, then checks all associativity equations directly.
The two levels have respectively two and $167$ raw assignments,
of which two and sixty-seven are associative.  Their entire
hom-set-permutation orbit sets agree with the completed streams.
Every one of these sixty-nine labelled tables is then subjected
to the full-basis cup checker.  This fresh check verifies $9280$
typed associativity equations and all eight ordered basis
products occurring before removal of repeated table orbits.
Every product has the stored explicit primitive.  The final
record \path{code/outputs/d_five_chain_independent.json} has
status \texttt{PASS}; the separate orbit-level record is
\path{code/outputs/d_five_chain_all_cups.json}.
Linearity of the cup product now proves the assertion for all
classes.  The argument uses only mod-two cohomology and makes
no inference about divisibility of integral cup products.
\end{proof}

\begin{theorem}[Acyclic support with trivial endomorphisms]
\label{C:all-acyclic21}
Let $C$ be a finite category with trivial endomorphism monoids
and acyclic support on distinct objects.  If $C$ has at most
twenty-one arrows, then $BC$ is not homotopy equivalent to
$\mathbb {CP}^2$.
\end{theorem}
\begin{proof}
We may assume that $C$ is connected.  Write $k$ for its number
of objects.  If its nerve has no four-simplex, fourth homology
vanishes.  Otherwise fix a four-simplex, and let $D$ be the
full subcategory on its five objects.  Those objects form a
linear order, and all ten forward hom-sets are nonempty.
If $k=5$, Proposition~\ref{C:five-object-acyclic21} applies;
if $k<5$, no four-simplex exists.  We treat the remaining
object counts by their cross-arrow budget.

Suppose a second four-simplex uses $r$ objects outside $D$.
For $1\leq r\leq4$, its five-object chain forces at least
\[
 r(5-r)+\binom r2
\]
new cross-hom sets involving those $r$ objects.  Its union
with the original chain is connected because they share an
object.  Connecting the other $k-5-r$ objects therefore
requires at least $k-5-r$ further cross arrows.  Including
the ten original cross arrows and the $k$ identities, the
total is at least
\begin{equation}\label{C:outside-four-simplex-cost}
 k+10+r(5-r)+\binom r2+(k-5-r)
 =2k+5+\frac{7r-r^2}{2}.
\end{equation}
For $r=1,2,3,4$ these bounds are respectively
$2k+8,2k+10,2k+11,2k+11$.  If $r=5$, the two chains are
disjoint.  The second chain contributes ten new cross arrows;
one further arrow is needed to connect the two components,
and the remaining $k-10$ objects need at least that many
additional arrows.  The total is then at least $2k+11$.
These estimates count only required nonempty hom-sets, so
parallel arrows can only increase them.

For $k\geq7$, all these bounds exceed twenty-one.  Hence
every four-simplex lies in $D$.  Every five-simplex would
have a four-dimensional face using an outside object, so
there are no five-simplices.  Thus the normalized chain
groups and boundary maps give an actual isomorphism
\[
 H_4(BD;\mathbb F_2)\longrightarrow H_4(BC;\mathbb F_2).
\]
Connecting the $k-5$ outside objects costs at least $k-5$
cross arrows as well as their identities.  Consequently
\[
 |D|\leq21-2(k-5)\leq17.
\]
The top-face argument in
Proposition~\ref{C:five-object-acyclic21} makes $H_4(BD)$
zero at this size, a contradiction.

It remains to consider $k=6$, with outside object $v$.
First suppose that $v$ lies in a four-simplex.  It is
comparable in the support order with at least four of the
five objects of $D$.  If it is comparable with all five,
the six objects form a chain.  The fifteen required cross
arrows and six identities exhaust the budget, so $C$ is
the ordinary six-element chain poset and is contractible.

Otherwise $v$ is incomparable with exactly one object $a$
of $D$.  Transitivity forces $v$ and $a$ to have the same
order position relative to the other four objects.  The
support is therefore a five-level order with one doubled
level.  It has fourteen required cross arrows and six
identities, leaving room for at most one additional
parallel arrow.  Its only five-object chains are $D$ and
the chain obtained by replacing $a$ with $v$.

The normalized four-chains consequently split into two
summands, one for each choice at the doubled level.  In
either summand, choose a face which deletes neither its
exclusive vertex nor either endpoint of the possible
parallel adjacent arrow.  At most three of the five
vertices are forbidden, so such a face exists.  It retains
the exclusive vertex, hence its target three-chain summand
receives no boundary contribution from the other branch.
If the parallel arrow occurs among the four adjacent arrows,
the selected face also retains its label; all other arrows
are unique.  This face map is therefore injective on the
four-simplex basis.  If the extra arrow is not adjacent in
this chain, that basis has one element and the same conclusion
holds.  Projecting the boundary to the two selected face
summands proves that $\partial_4$ is injective.  Again
$H_4(BC;\mathbb F_2)=0$.

Finally suppose that $v$ lies in no four-simplex.  All
four-simplices and all their faces are literally in $D$,
and there are no five-simplices.  Inclusion again induces
an isomorphism on fourth mod-two homology, hence restriction
induces an isomorphism
\[
 H^4(BC;\mathbb F_2)\longrightarrow H^4(BD;\mathbb F_2).
\]
The outside object requires its identity and at least one
cross arrow, so $|D|\leq19$.  By
Lemma~\ref{C:five-chain-nineteen-products}, every product
of degree-two classes on $D$ is zero.  Naturality and this
fourth-cohomology isomorphism imply that every degree-two
product on $C$ is zero as well.  This contradicts the
nonzero square of the generator of
$H^2(\mathbb {CP}^2;\mathbb F_2)$ and completes all cases.
\end{proof}
